%% Generated by tables/make_wood_tables.py -- do not edit.
\begin{deluxetable*}{lccccccccccc}
\tablecaption{Lexington \HII{} region benchmark HII20 (20~kK blackbody):
the published multi-code results of \citet[][their Table 2]{wood2004}, with a
current \textsc{Cloudy} c25.00 run and the \MoCHII{} result appended.  Line
intensities are relative to H$\beta$.\label{tab:wood20}}
\tablewidth{0pt}
\tabletypesize{\footnotesize}
\tablehead{\colhead{Quantity} & \colhead{GF} & \colhead{C25} & \colhead{HN} & \colhead{DP} & \colhead{TK} & \colhead{PH} & \colhead{RS} & \colhead{RR} & \colhead{BE} & \colhead{WME} & \colhead{\MoCHII{}}}
\startdata
$L(\mathrm{H}\beta)/10^{36}\,\mathrm{erg\,s^{-1}}$ & 4.85 & 4.99 & 4.85 & 4.83 & 4.90 & 4.93 & 5.04 & 4.89 & 4.97 & 4.87 & 4.71 \\
\HeI{} 5876 & 0.0072 & 0.0061 & 0.008 & 0.0073 & 0.008 & 0.0074 & 0.0110 & \nodata & 0.0069 & 0.00684 & 0.0068 \\
C\,\textsc{ii}] 2325+ & 0.054 & 0.076 & 0.047 & 0.046 & 0.040 & 0.060 & 0.038 & 0.063 & 0.042 & 0.053 & 0.068 \\
$[$N\,\textsc{ii}$]$ 122\,$\mu$m & 0.068 & 0.070 & \nodata & 0.072 & 0.007 & 0.072 & 0.071 & 0.071 & 0.071 & 0.067 & 0.073 \\
$[$N\,\textsc{ii}$]$ 6584+6548 & 0.745 & 0.809 & 0.786 & 0.785 & 0.925 & 0.843 & 0.803 & 0.915 & 0.846 & 0.778 & 0.840 \\
$[$N\,\textsc{ii}$]$ 5755 & 0.0028 & 0.0032 & 0.0024 & 0.0023 & 0.0029 & 0.0033 & 0.0030 & 0.0033 & 0.0025 & 0.0025 & 0.0032 \\
$[$N\,\textsc{iii}$]$ 57.3\,$\mu$m & 0.0040 & 0.0038 & 0.0030 & 0.0032 & 0.0047 & 0.0031 & 0.0020 & 0.0022 & 0.0019 & 0.0049 & 0.0039 \\
$[$O\,\textsc{i}$]$ 6300+6363 & 0.0080 & 0.0084 & 0.0060 & 0.0063 & 0.0059 & 0.0047 & 0.0050 & \nodata & 0.0088 & 0.055 & 0.0112 \\
$[$O\,\textsc{ii}$]$ 7320+7330 & 0.0087 & 0.0113 & 0.0085 & 0.0089 & 0.0037 & 0.0103 & 0.0080 & 0.0100 & 0.0064 & 0.0089 & 0.0122 \\
$[$O\,\textsc{ii}$]$ 3726+3729 & 1.01 & 1.40 & 1.13 & 1.10 & 1.04 & 1.22 & 1.08 & 1.17 & 0.909 & 1.12 & 1.41 \\
$[$O\,\textsc{iii}$]$ 52+88\,$\mu$m & 0.0030 & 0.0033 & 0.0026 & 0.0026 & 0.0040 & 0.0037 & 0.0020 & 0.0017 & 0.0022 & 0.0044 & 0.0033 \\
$[$O\,\textsc{iii}$]$ 5007+4959 & 0.0021 & 0.0023 & 0.0016 & 0.0015 & 0.0024 & 0.0014 & 0.0010 & 0.0010 & 0.0011 & 0.0026 & 0.0021 \\
$[$Ne\,\textsc{ii}$]$ 12.8\,$\mu$m & 0.264 & 0.296 & 0.267 & 0.276 & 0.27 & 0.271 & 0.286 & 0.290 & 0.295 & 0.289 & 0.303 \\
$[$S\,\textsc{ii}$]$ 6716+6731 & 0.499 & 0.744 & 0.473 & 0.459 & 1.02 & 0.555 & 0.435 & 0.492 & 0.486 & 0.521 & 0.778 \\
$[$S\,\textsc{ii}$]$ 4068+4076 & 0.022 & 0.023 & 0.017 & 0.020 & 0.052 & 0.017 & 0.012 & 0.015 & 0.013 & 0.017 & 0.023 \\
$[$S\,\textsc{iii}$]$ 18.7\,$\mu$m & 0.445 & 0.209 & 0.460 & 0.441 & 0.34 & 0.365 & 0.398 & 0.374 & 0.371 & 0.381 & 0.215 \\
$[$S\,\textsc{iii}$]$ 9532+9069 & 0.501 & 0.293 & 0.480 & 0.465 & 0.56 & 0.549 & 0.604 & 0.551 & 0.526 & 0.602 & 0.303 \\
$10^{3}\times\Delta(\mathrm{BC}\,3645)$/\AA & 5.54 & \nodata & \nodata & 5.62 & \nodata & 5.57 & 5.50 & \nodata & 6.18 & 5.52 & \nodata \\
$T_{\mathrm{inner}}$ / K & 7224 & 7155 & 6815 & 6789 & 6607 & 6742 & 6900 & 6708 & 6562 & 6852 & 6938 \\
$\langle T[N_{\mathrm{p}}N_{\mathrm{e}}]\rangle$ / K & 6680 & 6910 & 6650 & 6626 & 6662 & 6749 & 6663 & 6679 & 6402 & 6692 & 6916 \\
$R_{\mathrm{out}}/10^{18}$\,cm & 8.89 & 8.97 & 8.88 & 8.88 & 8.7 & 8.95 & 9.01 & 8.92 & 8.89 & 8.73 & 8.89 \\
$\langle\mathrm{He^{+}}\rangle/\langle\mathrm{H^{+}}\rangle$ & 0.048 & 0.043 & 0.051 & 0.049 & 0.048 & 0.044 & 0.077 & 0.034 & 0.041 & 0.047 & 0.046 \\
\enddata
\tablecomments{Codes, in the notation of \citet{wood2004}: GF =
Ferland's \textsc{Cloudy}; HN = Netzer's \textsc{Ion}; DP = P\'equignot's
\textsc{Nebu}; TK = Kallman's \textsc{XStar}; PH = Harrington's code;
RS = Sutherland's \textsc{Mappings}; RR = Rubin's \textsc{Nebula};
BE = Ercolano's \textsc{Mocassin}; WME = the code of \citet{wood2004}
itself.  Columns GF--BE agree entry by entry with Tables 4 and 5 of
\citet{ercolano2003}, the compilation of the published columns.
C25 is \textsc{Cloudy} c25.00 \citep{cloudy2025} run for this paper, and
is included because the published GF column dates from the 1990s version of
the same code, so the two together measure how far that reference has moved.
Its physics input is the unmodified c25.00 test-suite file
\code{tsuite/auto/hii\_paris.in} (HII40) or \code{hii\_coolstar.in} (HII20)
--- \textsc{Cloudy}'s own maintained version of these benchmarks --- with
only output commands appended; line luminosities come from a \code{save line
list absolute} file and are divided by the \textsc{Cloudy} H\,\textsc{i}
4861.32\,\AA{} luminosity, and the \textsc{Cloudy} blends summed for each row
are listed in \code{tables/make\_wood\_tables.py}.  Unlike the published
columns, C25 runs to the stopping criterion of its own deck rather than to a
prescribed outer radius --- an electron fraction of $10^{-2}$ at HII40, the
1000~K temperature floor at HII20 --- which mainly affects the
low-ionization rows fed by the partly neutral edge.
The \MoCHII{} column uses continuous stellar photon energies sampled from
the Planck spectrum with a Sobol sequence (seed 20260728), and evaluates
photoionization cross sections at each sampled energy; the 32-bin ionizing
array is a diagnostic tally only.  It uses the explicit Monte Carlo diffuse
field including the \HeI{} channels (case~A, matching the transfer treatment
of the published codes), and Badnell 2023 recombination with the Mao \&
Kaastra 2016 ground-level split, at $2^{25}$ stellar packets with the
volume-integrated convergence criterion; the ionizing luminosity is set from
the benchmark $L({\rm BB})$ so that $Q({\rm H})$ matches the specification;
blends are the sum of the
\MoCHII{} components listed in \code{tables/make\_wood\_tables.py}, and
\MoCHII{} wavelengths are vacuum values.  The \MoCHII{} structural entries
are computed from the same run: $T_{\mathrm{inner}}$ is the mean $\Te$ of the
cells at the illuminated inner edge, $R_{\mathrm{out}}$ the radius at which
the shell-averaged $x_{\mathrm{HII}}$ falls through 0.5, and the helium ratio
follows Eq.~(17) of \citet{ercolano2003},
$\langle\mathrm{He^{+}}\rangle/\langle\mathrm{H^{+}}\rangle =
\int \nelec n(\mathrm{He^{+}})\,{\rm d}V / \int \nelec n_{\mathrm{p}}\,{\rm d}V
\times n_{\mathrm{H}}/n_{\mathrm{He}}$.  Where the ionized gas still fills the
outermost shell of the prescribed sphere, $R_{\mathrm{out}}$ would be set by the
grid rather than predicted, and the entry is left empty.  The Balmer jump is
not computed.}
\end{deluxetable*}
